\section{Finite tests on residual curves}\label{sec:residual-tests}\label{sec:tests}

We next construct residual curves that retain any prescribed finite collection
of Kummer relations.  The curve can depend on the functions and on the
exponent.  Its genus, however, is bounded in terms of one fixed projective
model.  This distinction will make the pencil estimate uniform.

Fix a normal integral projective model $X\subset\mathbb P^N_l$ of $L/l$,
and put $d=\operatorname{trdeg}(L/l)\ge2$.
For a prime divisor $D$ on $X$, let $\deg D$ denote its degree in this
embedding.  For an integer $m\ge2$ and $g\in L^\times$, set
\[
 s_m(g)=\sum_{\substack{D\subset X\text{ prime divisor}\\
                       \operatorname{ord}_D(g)\not\equiv0\pmod m}}
                 \deg D.
\]
This is a finite sum.

\begin{lemma}[Finite residual tests]\label{tests:residual}
There is an integer $G_X\ge0$, depending only on the embedded variety $X$,
with the following property.  Given $g_1,\ldots,g_r\in L^\times$ and
$q=\ell^e$, $e\ge1$, there is a quasi-prime $(d-1)$-divisor $w$ of $L/l$
such that:
\begin{enumerate}
\item $lw$ is an algebraic closure of a finite field of characteristic
      different from $\ell$, and $Lw=lw(C)$ for a smooth projective integral
      curve $C/lw$ with $g(C)\le G_X$;
\item $w(g_j)=0$ for every $j$, and, for every
      $(n_1,\ldots,n_r)\in(\mathbb Z/q)^r$,
\[
 \prod_jg_j^{n_j}\in L^{\times q}
 \quad\Longleftrightarrow\quad
 \prod_j\overline g_j^{n_j}\in(Lw)^{\times q};
\]
\item for $m=\ell$ and $m=q$,
\begin{equation}\label{tests:support}
 \#\{P\in C(lw):\operatorname{ord}_P(\overline g_j)
                                  \not\equiv0\pmod m\}=s_m(g_j)
 \qquad(1\le j\le r).
\end{equation}
\end{enumerate}
Here $\overline g_j$ is the residue of $g_j$ at $w$.
\end{lemma}

We use the following form of Bertini irreducibility.  Its hypotheses include
the finite covering maps that occur below.

\begin{lemma}[Hyperplane sections of a quasi-finite map]\label{tests:bertini}
Let $V$ be an integral variety over an algebraically closed field, and let
$f:V\to\mathbb P^N$ be quasi-finite, with $\dim V=a\ge2$.
For a general hyperplane $H$, the inverse image $f^{-1}(H)$ is
geometrically irreducible.  If $f$ factors as a finite \'etale map to a
smooth locally closed subvariety $U\subset\mathbb P^N$, this inverse image
is also smooth for a general $H$.
\end{lemma}

\begin{proof}
We recall the incidence argument for irreducibility; see also
\cite[Lemma~37.32.3, Tag~0G4F]{Stacks}.
Write $B=(\mathbb P^N)^*$ and
\[
 I=\{(x,H)\in V\times B:f(x)\in H\}.
\]
The projection to $V$ is a projective bundle, so $I$ is integral.
In $I\times_BI$, the open set lying over pairs $(x,y)$ with
$f(x)\ne f(y)$ is a projective bundle with fiber $\mathbb P^{N-2}$ over
an irreducible open subset of $V\times V$.  It is therefore irreducible
of dimension $2a+N-2$.
The locus of pairs with $f(x)=f(y)$ has dimension at most $a$, since $f$
is quasi-finite; the incidence above it has dimension at most $a+N-1$.
This is strictly smaller than $2a+N-2$.

Every component of the generic hyperplane section has dimension $a-1$:
a hyperplane cuts a nonzero nonunit in an integral variety of finite type
over a field.  Consequently every component of its square has dimension
$2a-2$.  The equal-image locus cannot supply a component of this generic
square.  The generic square is thus irreducible.  Equivalently, the
function field of $B$ is separably algebraically closed in that of $I$,
so the generic hyperplane section is geometrically irreducible.
Geometric irreducibility spreads to a nonempty open of $B$
\cite[Lemma~37.27.5, Tag~0559]{Stacks}.
Finally, a general hyperplane section of the smooth embedded variety $U$
is smooth by Bertini, and its \'etale inverse image is smooth.
\end{proof}

\begin{proof}[Proof of Lemma~\ref{tests:residual}]
\emph{Spread the finite data.}
Let $D_1,\ldots,D_s$ be the prime divisors occurring in the divisors of
the $g_j$.  There is a closed subset $B_X\subset X$ of codimension at
least two such that $X\setminus B_X$ is smooth, the
$D_i\setminus B_X$ are smooth pairwise disjoint Cartier divisors, and,
near each of them, every $g_j$ is a power of a local equation times a
unit.  Off their union the $g_j$ are units.
These assertions follow by deleting the singular loci, the pairwise
divisor intersections and the proper closed exceptional subsets on
which the required local identities fail.

Choose a smooth dense open $U\subset X$ on which all $g_j$ are units.
The simultaneous root cover
\[
 V=\operatorname{Spec}_U
 \mathcal O_U[T_1,\ldots,T_r]/(T_1^q-g_1,\ldots,T_r^q-g_r)
 \longrightarrow U
\]
is finite \'etale of degree $q^r$.
Let $H$ be the subgroup generated by the $[g_j]$ in
$L^\times/L^{\times q}$.  Kummer theory shows that each connected
component of $V$ has degree $|H|$ over $U$.
Since $U$ is smooth, these components are integral.

All these data descend to a finitely generated domain $R\subset l$ in
which $\ell$ is invertible.  Enlarge $R$ and then localize it so that the
following properties hold in every geometric fiber:
\begin{itemize}
\item the embedded model $X_R$ is projective and flat, with normal
      integral fibers and the Hilbert polynomial of $X$;
\item the $D_{i,R}$ are distinct geometrically integral divisors with
      their original degrees, and the local divisor identities and
      smoothness conditions above hold outside a subset of fiberwise
      codimension at least two;
\item the root cover and its disjoint decomposition into components
      extend, and every component has geometrically integral fibers
      and finite locally free degree $|H|$ over $U_R$.
\end{itemize}
These are finite spreading conditions.  Flatness makes the projective
Hilbert polynomials constant.  Geometric integrality and normality
spread from the generic fiber after shrinking; the latter uses the
proper flat finite-presentation form of
\cite[Theorem~12.2.4(iv),(viii)]{EGAIV3}.
The equations $g_j=u\pi^n$, the inverses specifying units, the component
idempotents, and the finite locally free ranks are all finite algebraic
data.  Spreading them preserves the exact orders and degrees, rather
than just the supports.  The dimension bound on the omitted closed set
also persists after shrinking.

For later use include a separating transcendence basis
$x_1,\ldots,x_d$ of $L/l$, a primitive element $y$ over $l(x_1,\ldots,x_d)$,
its monic minimal polynomial, and common birational opens with $X$.
After another localization, the same presentation and birational
identification hold in the geometric fibers, with the reduced polynomial
irreducible of its original degree over the rational function field.
Preserve as well the nonzero reductions of the finitely many
numerators and denominators used for the $g_j$.
This follows either by spreading those common opens and geometric
integrality, or by spreading the absolutely irreducible defining
polynomial after clearing denominators.

\emph{Realize a good fiber as a residue field.}
Choose a closed point $s\in\operatorname{Spec}R$.
Its residue field is finite and has characteristic different from $\ell$.
A valuation ring of $l$ dominating $R_s$ exists
\cite[Lemma~10.50.2, Tag~00IA]{Stacks}.
Its residue field is algebraically closed.  If that residue field is
not algebraic over a finite field, give a transcendence basis independent
values in an ordered free abelian group, prolong the resulting valuation
to the residue field, and compose.  This yields a valuation $s_0$ of $l$
with the same center $s$ and residue field $\lambda$ algebraic over a
finite field.  Residues of algebraically closed fields are algebraically
closed, so $\lambda$ is an algebraic closure of a finite field.
The composition does not change the center because every nonzero
element of a finite field has value zero.
This construction uses a transcendence basis of arbitrary cardinality;
there is no restriction on the size of $l$.

Extend $s_0$ to $l(x_1,\ldots,x_d)$ by the Gauss valuation, whose residue
field is $\lambda(x_1,\ldots,x_d)$ and whose value group is $s_0l$.
Prolong it to $L$.  The element $y$ is integral and its residue satisfies
the reduced irreducible polynomial, of degree
$[L:l(x_1,\ldots,x_d)]$.  The fundamental valuation inequality therefore
forces the residue degree to be that degree and the value-group index
to be one \cite[Section~15.125, Tag~0ASF]{Stacks}.
We obtain a valuation $s_X$ with
\[
 s_XL=s_0l,\qquad Ls_X=\lambda(X_s),
\]
and with the $g_j$ reducing to their specified rational functions on
$X_s$.  This also explains why no extra ramification or defect is hidden
in the passage to the good fiber.

\emph{Cut down to a curve.}
Over $\lambda$, choose a general flag of $d-1$ hyperplanes in the fixed
projective embedding.  Make the final section $C$ avoid the omitted
codimension-two set and meet the $D_{i,s}$ transversely.
The intersections with different $D_{i,s}$ are disjoint.
The final section is smooth and integral.  Apply
Lemma~\ref{tests:bertini} successively to the finite list of covering
components as well as to the reduced underlying base sections; each restricted covering
component remains integral over $U_s\cap C$.
At every step the varieties to which irreducibility is applied have
dimension at least two, including the last surface-to-curve step.
Smoothness of the restricted cover follows from its \'etaleness over the
smooth curve open.

For completeness, intermediate scheme sections need only have
irreducible support and be smooth at the generic points used by the
flag.  All those points lie in the original smooth locus.
Choose the hyperplanes also to avoid the finitely many associated
points at each stage, so their equations are non-zero-divisors.
The hyperplane exact sequences give
\[
 P_C(n)=\Delta^{d-1}P_X(n),\qquad
 \Delta P(n)=P(n)-P(n-1).
\]
As $C$ is smooth integral, its genus is consequently fixed by the
embedded model $X$; denote it by $G_X$.
All conditions imposed on the flag are finitely many nonempty open
conditions.  They can be met over the infinite algebraically closed
field $\lambda$, even when it is countable.

Compose $s_X$ with the successive divisorial valuations of this flag.
For the first composition, the relative value group is the innermost
convex copy of $\mathbb Z$, and its residue transcendence degree drops
by one.  Every proper coarsening kills this copy of $\mathbb Z$ and
therefore has all its values supplied by constants.
The first composition is thus a quasi-prime divisor.  The remaining
flag valuations are prime divisors over the residue constants, giving
a quasi-prime $(d-1)$-divisor $w$.
All $g_j$ are units at every generic point in the flag, so $w(g_j)=0$,
and $Lw=\lambda(C)$.

\emph{Check relations and orders.}
An old relation $\prod_jg_j^{n_j}=b^q$ specializes: since $w(g_j)=0$,
torsion-freeness of the value group gives $w(b)=0$.
Thus the old relation kernel in $(\mathbb Z/q)^r$ is contained in the
new one.
On the other hand, every component of the root cover restricted to
$U_s\cap C$ is integral and has degree $|H|$.
Kummer theory gives the same order $|H|$ for the subgroup generated by
the residue classes.  The two finite relation kernels therefore have
the same order, so they are equal.

Finally, $C$ meets $D_{i,s}$ in exactly $\deg D_i$ distinct transverse
points.  At each such point the order of $\overline g_j$ equals the
unchanged integer $\operatorname{ord}_{D_i}(g_j)$.
There are no other zeros or poles.  This proves
\eqref{tests:support}, in fact for every divisor $m$ of $q$, and
completes the construction.
\end{proof}
